% Part: sets-functions-relations
% Chapter: sets
% Section: unions-and-intersections

\documentclass[../../../include/open-logic-section]{subfiles}

\begin{document}

\olfileid{sfr}{set}{uni}
\olsection{Verenigingen en doorsneden}

\begin{explain}
In \olref[sfr][set][bas]{sec} hebben we verzamelingen door abstractie
gedefinieerd, dus met definities van de vorm $\Setabs{x}{\phi(x)}$.
Daarbij gebruiken we een eigenschap~$\phi$, waarin ook al gedefinieerde
verzamelingen mogen voorkomen. Zijn $A$ en~$B$ bijvoorbeeld verzamelingen,
dan bestaat $\Setabs{x}{x \in A \lor x \in B}$ uit alle objecten die
!!{element} zijn van $A$ of van~$B$. Deze verzameling verenigt dus de
!!{element}en van $A$ en~$B$. In \olref{fig:union} geeft het gemarkeerde
gebied de !!{element}en van de twee verzamelingen $A$ en~$B$ aan.

\begin{figure}
  \olasset{assets/diagrams/union.tikz}
  \caption{De vereniging $A \cup B$ van twee verzamelingen bestaat uit de
   !!{element}en van $A$ samen met die van~$B$.}
  \ollabel{fig:union} 
\end{figure}

Verzamelingen op deze manier verenigen is een veelgebruikte en nuttige
bewerking. Daarom geven we haar een formele naam en een symbool.
\end{explain}

\begin{defn}[Vereniging]
De \emph{vereniging} van twee verzamelingen $A$ en $B$, geschreven als
$A \cup B$, is de verzameling van alle objecten die !!{element} zijn van
$A$, van $B$ of van beide.
\[
A \cup B = \Setabs{x}{x \in A \lor x \in B}
\]
\end{defn}

\begin{ex}
Omdat herhaling van !!{element}en geen rol speelt, bevat de vereniging van
twee verzamelingen een gemeenschappelijk !!{element} maar één keer. Zo geldt
$\{ a, b, c\} \cup \{ a, 0, 1\} = \{a, b, c, 0, 1\}$.

De vereniging van een verzameling met een van haar deelverzamelingen is
de grotere verzameling: $\{a, b, c \} \cup \{a \} = \{a, b, c\}$.

De vereniging van een verzameling met de lege verzameling is die verzameling
zelf: $\{a, b, c \} \cup \emptyset = \{a, b, c \}$.
\end{ex}

\begin{prob}
Bewijs dat uit $A \subseteq B$ volgt dat $A \cup B = B$.
\end{prob}

\begin{explain}
We kunnen ook de ``duale'' bewerking van vereniging beschouwen. Deze bewerking
vormt de verzameling van alle !!{element}en die zowel !!{element} van~$A$
als !!{element} van~$B$ zijn. Zij heet \emph{doorsnede} en kan worden
weergegeven zoals in \olref{fig:intersection}.
\begin{figure}
  \olasset{assets/diagrams/intersection.tikz}
  \caption{De doorsnede $A \cap B$ van twee verzamelingen bestaat uit hun
    gemeenschappelijke !!{element}en.}
  \ollabel{fig:intersection}
\end{figure}
\end{explain}

\begin{defn}[Doorsnede]
De \emph{doorsnede} van twee verzamelingen $A$ en $B$, geschreven als
$A \cap B$, is de verzameling van alle objecten die zowel !!{element} van
$A$ als !!{element} van~$B$ zijn.
\[
A \cap B = \Setabs{x}{x \in A \land x \in B}
\]
Twee verzamelingen heten \emph{disjunct} als hun doorsnede leeg is. Ze hebben
dan geen !!{element}en gemeen.
\end{defn}

\begin{ex}
Als twee verzamelingen geen !!{element}en gemeen hebben, is hun doorsnede leeg:
$\{ a, b, c\} \cap \{ 0, 1\} = \emptyset$.

Als twee verzamelingen wel !!{element}en gemeen hebben, bestaat hun doorsnede
uit al die !!{element}en: $\{a, b, c \} \cap \{a, b, d \} = \{a, b\}$.

De doorsnede van een verzameling met een van haar deelverzamelingen is de
kleinere verzameling: $\{a, b, c\} \cap \{a, b\} = \{a, b\}$.

De doorsnede van elke verzameling met de lege verzameling is leeg: $\{a, b, c \}
\cap \emptyset = \emptyset$.
\end{ex}

\begin{prob}
Geef een rigoureus bewijs dat uit $A \subseteq B$ volgt dat $A \cap B = A$.
\end{prob}

\begin{explain}
We kunnen ook de vereniging of doorsnede van meer dan twee verzamelingen
vormen. Een elegante algemene aanpak is de volgende. Breng alle verzamelingen
waarvan we de vereniging of doorsnede willen nemen onder in één verzameling.
De vereniging van de oorspronkelijke verzamelingen is dan de verzameling van
alle objecten die tot ten minste één !!{element} van die verzameling behoren;
hun doorsnede is de verzameling van alle objecten die tot elk !!{element}
ervan behoren.
\end{explain}

\begin{defn}
Als $A$ een verzameling van verzamelingen is, dan is $\bigcup A$ de verzameling
van de !!{element}en van de !!{element}en van~$A$:
\begin{align*}
\bigcup A & = \Setabs{x}{x \text{ behoort tot !!a{element} van } A},
\text{ d.w.z.:}\\
& = \Setabs{x}{\text{er is een } B \in A
  \text{ waarvoor } x \in B}
\end{align*}
\end{defn}

\begin{defn}
Als $A$ een verzameling van verzamelingen is, dan is $\bigcap A$ de verzameling
van de objecten die alle elementen van~$A$ gemeen hebben:
\begin{align*}
\bigcap A & = \Setabs{x}{x \text{ behoort tot elk !!{element} van } A},
\text{ d.w.z.:}\\
 & = \Setabs{x}{\text{voor elke } B \in A, x \in B}
\end{align*}
\end{defn}

\begin{ex}
Stel $A = \{ \{ a, b \}, \{ a, d, e \}, \{ a, d \} \}$.
Dan is $\bigcup A = \{ a, b, d, e \}$ en $\bigcap A = \{ a \}$.
\end{ex}
\begin{prob}
	Toon aan dat als $A$ een verzameling is en $A \in B$, dan
	$A \subseteq \bigcup B$.
\end{prob}

Hetzelfde kan voor een rij verzamelingen $A_1$, $A_2$, \dots
\begin{align*}
\bigcup_i A_i & = \Setabs{x}{x \text{ behoort tot ten minste één van de } A_i}\\
\bigcap_i A_i & = \Setabs{x}{x \text{ behoort tot elk van de } A_i}.
\end{align*}

Als we een \emph{indexverzameling} $I$ hebben en verzamelingen $A_i$
beschouwen, één voor elke $i \in I$, kunnen we ook deze afkortingen gebruiken:
\begin{align*}
	\bigcup_{i \in I} A_i & = \bigcup \Setabs{A_i }{i \in I}\\
	\bigcap_{i \in I} A_i & = \bigcap\Setabs{A_i}{i \in I}
\end{align*}

Ten slotte kunnen we de verzameling van alle !!{element}en van~$A$ beschouwen
die niet in~$B$ liggen. Dit is weergegeven in \olref{difference}.

\begin{figure}
  \olasset{assets/diagrams/difference.tikz}
  \caption{Het verschil $A \setminus B$ van twee verzamelingen bestaat uit
    de !!{element}en van~$A$ die geen !!{element} van~$B$ zijn.}
  \ollabel{difference}
\end{figure}

\begin{defn}[Verschil]
Het \emph{verschil}~$A \setminus B$ is de verzameling van alle !!{element}en
van $A$ die geen !!{element} van~$B$ zijn, dus
\[
A\setminus B = \Setabs{x}{x\in A \text{ en } x \notin B}.
\]
\end{defn}

\begin{prob}
	Bewijs dat uit $A \subsetneq B$ volgt dat $B \setminus A \neq \emptyset$.
\end{prob}
\end{document}
